% ============================================================================
%  SECCIÓN 5.7: EVIDENCIAS DEL AVANCE 5
% ============================================================================

\section{Evidencias del avance}

A continuaci\'on se presentan las evidencias del avance correspondiente al
PDF~5, organizadas seg\'un la gu\'ia: CRUD implementado, persistencia de
datos, capturas de pantalla de las mejoras y verificaci\'on tras cierre.

% --- EVIDENCIA 1 ---
\subsubsection*{Evidencia 1: Pantalla de gesti\'on de escenarios}

La nueva pantalla de gesti\'on (\textit{manage.tsx}) permite a los
usuarios con rol de administrador y gestor crear, editar y eliminar
escenarios deportivos. La lista muestra todos los escenarios (activos
e inactivos) con opciones de acci\'on:

\begin{figure}[H]
  \centering
  \caption{Evidencia: Pantalla de gesti\'on de escenarios}
  \label{fig:evidencia5-gestion}
  \contenidoConFuente{%
    \includegraphics[width=0.28\linewidth]{evidencia5-pantalla-gestion.png}%
  }{Elaboraci\'on propia (captura de la aplicaci\'on).}
\end{figure}

% --- EVIDENCIA 2 ---
\subsubsection*{Evidencia 2: Formulario de creaci\'on de escenario}

El formulario de creaci\'on (\textit{scenario-form/\text{[}id\text{]}.tsx})
incluye validaciones con Zod, obtenci\'on de ubicaci\'on GPS y subida
de im\'agenes. Los campos se validan en tiempo real antes de enviar
los datos a Supabase:

\begin{figure}[H]
  \centering
  \caption{Evidencia: Formulario de creaci\'on de escenario}
  \label{fig:evidencia5-formulario-crear}
  \contenidoConFuente{%
    \includegraphics[width=0.28\linewidth]{evidencia5-formulario-crear.png}%
  }{Elaboraci\'on propia (captura de la aplicaci\'on).}
\end{figure}

% --- EVIDENCIA 3 ---
\subsubsection*{Evidencia 3: Formulario de edici\'on de escenario}

El mismo formulario se reutiliza para editar un escenario existente,
cargando los datos actuales en los campos del formulario. Al enviar,
se ejecuta un \textit{upsert} que actualiza el registro en Supabase:

\begin{figure}[H]
  \centering
  \caption{Evidencia: Formulario de edici\'on de escenario}
  \label{fig:evidencia5-formulario-editar}
  \contenidoConFuente{%
    \includegraphics[width=0.28\linewidth]{evidencia5-formulario-editar.png}%
  }{Elaboraci\'on propia (captura de la aplicaci\'on).}
\end{figure}

% --- EVIDENCIA 4 ---
\subsubsection*{Evidencia 4: Eliminaci\'on de escenario (soft delete)}

La operaci\'on de eliminaci\'on suave cambia el estado del escenario
a ``inactivo'', removi\'endolo del cat\'alogo p\'ublico sin eliminarlo
de la base de datos. Solo los administradores pueden ejecutar esta
operaci\'on:

\begin{figure}[H]
  \centering
  \caption{Evidencia: Eliminaci\'on de escenario (soft delete)}
  \label{fig:evidencia5-eliminar}
  \contenidoConFuente{%
    \includegraphics[width=0.28\linewidth]{evidencia5-eliminar-escenario.png}%
  }{Elaboraci\'on propia (captura de la aplicaci\'on).}
\end{figure}

% --- EVIDENCIA 5 ---
\subsubsection*{Evidencia 5: Detalle de escenario con eventos}

La pantalla de detalle muestra la informaci\'on completa del
escenario, incluyendo las im\'agenes almacenadas en Supabase Storage,
los deportes disponibles y los eventos pr\'oximos. Los usuarios con
permiso pueden crear y eliminar eventos:

\begin{figure}[H]
  \centering
  \caption{Evidencia: Detalle de escenario con eventos}
  \label{fig:evidencia5-detalle}
  \contenidoConFuente{%
    
\includegraphics[width=0.28\linewidth]{evidencia5-detalle-eventos.png}%
  }{Elaboraci\'on propia (captura de la aplicaci\'on).}
\end{figure}

% --- EVIDENCIA 6 ---
\subsubsection*{Evidencia 6: Creaci\'on de evento}

El formulario de creaci\'on de evento permite a los gestores y
administradores agregar nuevas actividades a un escenario. Se
validan campos como nombre, fecha, hora y descripci\'on:

\begin{figure}[H]
  \centering
  \caption{Evidencia: Creaci\'on de evento}
  \label{fig:evidencia5-crear-evento}
  \contenidoConFuente{%
    \includegraphics[width=0.28\linewidth]{evidencia5-crear-evento.png}%
  }{Elaboraci\'on propia (captura de la aplicaci\'on).}
\end{figure}

% --- EVIDENCIA 7 ---
\subsubsection*{Evidencia 7: Eliminaci\'on de evento}

Los eventos pueden ser eliminados por los gestores y administradores.
La eliminaci\'on es inmediata y se refleja en la interfaz tras
invalidar el cach\'e de React Query:

\begin{figure}[H]
  \centering
  \caption{Evidencia: Eliminaci\'on de evento}
  \label{fig:evidencia5-eliminar-evento}
  \contenidoConFuente{%
    \includegraphics[width=0.28\linewidth]{evidencia5-eliminar-evento.png}%
  }{Elaboraci\'on propia (captura de la aplicaci\'on).}
\end{figure}

% --- EVIDENCIA 8 ---
\subsubsection*{Evidencia 8: Persistencia de sesi\'on tras cierre}

Para verificar la persistencia, se cerr\'o la aplicaci\'on completamente
y se reabri\'o. La sesi\'on del usuario se mantuvo activa sin solicitar
credenciales, demostrando que AsyncStorage conserva los tokens JWT:

\begin{figure}[H]
  \centering
  \caption{Evidencia: Sesi\'on persistente tras cierre}
  \label{fig:evidencia5-sesion}
  \contenidoConFuente{%
    \includegraphics[width=0.28\linewidth]{appmovil-ejecutando-fisico.jpeg}%
  }{Elaboraci\'on propia (captura de la aplicaci\'on).}
\end{figure}

% --- EVIDENCIA 9 ---
\subsubsection*{Evidencia 9: Favoritos persistentes}

Los favoritos guardados por el usuario persisten en Supabase y se
recuperan autom\'aticamente al reabrir la aplicaci\'on. La lista de
favoritos muestra los escenarios guardados con sus im\'agenes:

\begin{figure}[H]
  \centering
  \caption{Evidencia: Favoritos persistentes tras cierre}
  \label{fig:evidencia5-favoritos}
  \contenidoConFuente{%
    \includegraphics[width=0.28\linewidth]{evidencia-favorito-despues.jpg}%
  }{Elaboraci\'on propia (captura de la aplicaci\'on).}
\end{figure}

% --- EVIDENCIA 10 ---
\subsubsection*{Evidencia 10: Repositorio actualizado}

El historial de commits documenta la evoluci\'on del proyecto con las
nuevas funcionalidades de persistencia:

\begin{lstlisting}[style=terminal]
feat(manage): add scenario management screen
feat(scenario-form): add create/edit scenario form
feat(events): add event creation and deletion
feat(persistence): implement AsyncStorage for session
feat(supabase): add RLS policies and performance indexes
fix(storage): restrict image upload to admin role
\end{lstlisting}

La captura general del repositorio se presenta en la
figura~\ref{fig:repositorio-github} de la configuraci\'on de Git
y GitHub.
